\begin{frame}{Introduction to Superfluidity}
\begin{itemize}
    \item The classical fluid dynamics has been known for more than a century, however, a complete understanding of the model (expecially under turbulent regimes) is still missing \cite{frisch1995turbulence}
    \item In the context of quantum mechanics, the bonsons can occupy simoultaneously the same quantum state, leading to the formation of a superfluid phase \cite{pitaevskii2016bose}, characterized by lack of viscosity and quantized vorticity
    \item The study of superfluid turbulence can help to understand classical turbulence, since the quantization of vorticity simplifies the dynamics 
\end{itemize}
\end{frame}

\begin{frame}{The Gross-Pitaevskii Equation}
   The governing equation for a weakly interacting Bose-Einstein condendate is the \emph{Gross-Pitaevskii equation} \cite{pitaevskii2016bose}:
    \begin{equation*}
        \ii \hbar \, \partial_t \, \psi = -\frac{\hbar^2}{2m} \nabla^2 \psi  + V_0 |\psi|^2 \psi - E_0 \psi
    \end{equation*}

    where $\psi(\mathbf{r},t)$ is the macroscopic wave function of the condensate, $m$ is the mass of the bosons, $V_0$ is the interaction strength and $E_0$ is the chemical potential. It can be rewritten in dimensionless form  \cite{CZ12} as:
    \begin{equation*}
        \partial_t \, \psi = \frac{\ii}{2} \nabla^2 \psi + \frac{\ii}{2} (1 - |\psi|^2) \psi
    \end{equation*}
\end{frame}

\begin{frame}[allowframebreaks]{Classical analogy via Madelung transformation}
    To achive a classical analogy, we can use the \emph{Madelung transformation} \cite{madelung1927quantentheorie}:
    \begin{equation}
        \psi(\rr, t) = \rho^{1/2}\exp(\ii S)
        \label{eq:madelung}
    \end{equation}    

    where $\rho = \rho(\rr, t) = \abs{\psi(\rr, t)}^2$ is the density and $S = S(\rr, t)$ is the phase of $\psi$. The vector field $\uu = \nabla S$ can be interpreted as the superfluid velocity field.

    Substituting \eqref{eq:madelung} into the GPE and separating real and imaginary parts, we obtain the following system, written in tensor notation:

    \begin{align*}
        &\pdv{\rho}{t} + \pdv{\rho u_j}{x_j} = 0 \\
        &\rho \left(\pdv{u_i}{t} + u_j \pdv{u_i}{x_j}\right) = -\pdv{p}{x_i} + \pdv{\tau_{ij}}{x_j}, \quad i = 1,2,3
    \end{align*}

    We can observe that this system is very similar to the Navier-Stokes equations, except for the pressure and stress term. This to justify the introduction.
\end{frame}

\begin{frame}[allowframebreaks]{Approximation of steady-state vortex}
    In a two dimensional domain, we set $\psi(x,y) = f(r)\exp(\ii \theta(x,y)) \quad r = \sqrt{x^2 + y^2}$, where we imposed the time independency.
    In order to obtain an initial approximaion of the straight vortex, we substitute the latter equation into GPE, imposing the time independency. We obtain the boundary value problem

    \begin{equation*}
        f'' + \frac{f'}{r} + f\left(1 - f^2 - \frac{1}{r^2}\right) = 0
    \end{equation*}

    with boundary values $f(0) = 0$, $f(\infty) = 1$.

    One can integrate this equation numerically, however, it is convenient to approximate $f$ with a Padé approximation \cite{Berloff20041617} in order to evaluate the latter everywhere. The approximant used in what follows is a 4-th order diagonal Padé approximaion, which coefficents are reported in \cite{CZ18}.
\end{frame}